\scanpage{32}
% sheet: reused_typed
% hand: typed
% archive sheet number: 31
% Typescript rotated 180 degrees relative to the archive foliation.
% @ handwritten / top centre in reading orientation
\handwritten{\mspage{55 \unsure{ter|der}}}
% @ text / proof, opening
\uline{Démonstration}. Le cas a) n'est autre que 5.10 et a été mis pour mémoire.
% @ text / proof, part b
b) \deleted{Le cas $\ell\neq p$ n'est autre que le critère galoisien de \unsure{la|ré}} \inserted{La première assertion est une conséquence \deleted{immé} triviale de} a) et des définitions. Pour prouver la dernière assertion dans b), on peut \deleted{donc} supposer que $\ell\neq p$, mais alors grâce à a) il suffit de prendre pour $K'$ une extension finie séparable de $K$ correspondant à un sous-groupe ouvert $\handwritten{\pi'}$ de $\handwritten{\pi}=\mathrm{Gal}(\bar K/K)$ dont l'intersection avec le groupe d'inertie $I$ est contenue dans le \deleted{stabli} noyau de la représentation de $I$ sur $T_{\ell}(A_K(\bar K))$, lequel est \deleted{\unsure{fi|à}} ouvert par hypothèse (cf. 5.12.1).
% @ text / proof, first three implications
c) On a déjà signalé que $1)\handwritten{\Rightarrow}2)$ (5.12.2), il est trivial que $2)\handwritten{\Rightarrow}3)$, reste à prouver que $3)\handwritten{\Rightarrow}1)$. Pour $\ell\neq p$, ce n'est autre que \inserted{le} cri\deleted{\unsure{t|i}}tère galoisien de réduction semi-stable 3.5. Pour $\ell=p$, on peut procéder ainsi. \deleted{\unsure{A|O}}Choisissons un nombre premier $\ell'\neq p$, et appliquons à celui-ci les considérations d e 4.1. \deleted{On est donc ramené à prouver que la représentation}
% @ text / proof, strictly local case through end
Supposons $S$ strictement local, ce qui est loisible (3.4.0), prenons une extension galoisienne finie $K'$ de $K$, d\deleted{d}e groupe $G$, telle que $A_{K'}$ ait réduction semi-stable. Il faut prouver que la représentation de $G$ sur les trois quotients intervenant dans (4.1.1) est triviale, ou ce qui revient au même, que la \deleted{somme des traces} représentation somme est triviale, ou encore, que la trace de cette dernière est "triviale", i.e. constante . Or nous avons vu dans 4.3 a) que cette fonction trace est indépendante d\deleted{\unsure{a}}u choix de $\ell'\neq p$, et dans 6.\deleted{\unsure{6|5}}\handwritten{4} plus bas nous indiquons une méthode de calcul de cette fonction en \deleted{fonction}\inserted{termes} de $T_p(A_K)$ également. Ce calcul donne trivialement une fonction constante dans le cas où $T_p(A_K)$ a bonne réduction virtuelle, et cela achève la démonstration.
% @ text / bottom, remark, with inserted qualifier and manuscript postscript
\uline{Remarque} 5.13.1. Il semble plausible que si $X_{\handwritten{\eta}}$ est un \handwritten{\inserted{pro-$\unsure{l|\ell}$-}}groupe de Barsotti-Tate sur $\handwritten{\eta}$ qui a bonne réduction virtuelle sur $S$, et bonne réduction virtuelle essentielle d'échelon $n$, alors il a \deleted{\unsure{v|b}}bonne réduction virtuelle d'échelon $n$. \handwritten{(C'est \deleted{ce qui avait été utilisé, pour} vrai pour $l\neq p$.)}
